\section{Proofs}
\label{app:proofs}

\subsection{Proof of Proposition~1 (Hard-Sample Bias)}
\label{app:proof-hsb}

\textbf{Setup.}
Let $\mathcal{I}_\text{self}(z) = \nabla^\top H_{\hat\theta}^{-1} \nabla$, where $\nabla \equiv \nabla_\theta \ell(z;\hat\theta)$ and $H_{\hat\theta}$ is the empirical Hessian, assumed symmetric positive definite (SPD) with $\lambda_{\min}(H_{\hat\theta}) > 0$ (this holds whenever the model is not at a degenerate critical point).
Partition $\mathcal{D}$ into $\mathcal{D}_\text{clean}$ and $\mathcal{D}_\text{noisy}$ such that $\|\nabla\| \leq \epsilon$ for all $z \in \mathcal{D}_\text{clean}$ and $\|\nabla\| \geq \delta$ for all $z \in \mathcal{D}_\text{noisy}$, with $\delta \gg \epsilon$.

\textbf{Step~1 --- Bound $\mathcal{I}_\text{self}$ for each partition.}
Since $H_{\hat\theta}$ is SPD with eigenvalues in $[\lambda_{\min}(H_{\hat\theta}),\, \lambda_{\max}(H_{\hat\theta})]$, the Rayleigh quotient gives:
\begin{equation}
    \frac{\|\nabla\|^2}{\lambda_{\max}(H_{\hat\theta})}
    \;\leq\; \mathcal{I}_\text{self}(z)
    \;\leq\; \frac{\|\nabla\|^2}{\lambda_{\min}(H_{\hat\theta})}.
\end{equation}
Substituting the gradient-norm assumptions ($\|\nabla\|^2 \leq \epsilon^2$ for clean samples; $\|\nabla\|^2 \geq \delta^2$ for noisy samples) and taking expectations (the bounds are uniform over each partition, so expectation preserves them):
\begin{align}
    \mathbb{E}_{z \sim \mathcal{D}_\text{clean}}[\mathcal{I}_\text{self}(z)]
        &\leq \frac{\epsilon^2}{\lambda_{\min}(H_{\hat\theta})}, \\
    \mathbb{E}_{z \sim \mathcal{D}_\text{noisy}}[\mathcal{I}_\text{self}(z)]
        &\geq \frac{\delta^2}{\lambda_{\max}(H_{\hat\theta})}.
\end{align}

\textbf{Step~2 --- Lower bound on $\Delta_\text{HSB}$.}
\begin{equation}
    \Delta_\text{HSB}
    := \mathbb{E}_{z \sim \mathcal{D}_\text{noisy}}[\mathcal{I}_\text{self}(z;\hat\theta)]
     - \mathbb{E}_{z \sim \mathcal{D}_\text{clean}}[\mathcal{I}_\text{self}(z;\hat\theta)]
    \;\geq\; \frac{\delta^2}{\lambda_{\max}(H_{\hat\theta})} - \frac{\epsilon^2}{\lambda_{\min}(H_{\hat\theta})}.
\end{equation}
Substituting $\lambda_{\min}(H_{\hat\theta}) = \lambda_{\max}(H_{\hat\theta})/\mathrm{cond}(H_{\hat\theta})$:
\begin{equation}
    \Delta_\text{HSB}
    \;\geq\; \frac{\delta^2 - \mathrm{cond}(H_{\hat\theta})\,\epsilon^2}{\lambda_{\max}(H_{\hat\theta})}
    \;>\; 0
    \qquad \text{whenever} \quad \frac{\delta}{\epsilon} > \sqrt{\mathrm{cond}(H_{\hat\theta})}.
\end{equation}
Under the interpolation regime (an overparameterised model trained to near-zero training loss), $\epsilon \to 0$: the bound simplifies to $\delta^2/\lambda_{\max}(H_{\hat\theta}) > 0$ and the positivity condition is trivially satisfied ($\delta/\epsilon \to \infty$). $\square$

% ------------------------------------------------------------
\subsection{Proof of Proposition~2 (Trajectory Integration Reduces HSB)}
\label{app:proof-traj}
% ------------------------------------------------------------

\textbf{Proposition~2.}
Under smooth loss and bounded gradient noise, the trajectory-integrated knowledge score $\kappa_n$ produces a strictly smaller noisy-over-clean gap than the single last-checkpoint baseline $\kappa^\dagger_n = 1/I_{(n,T-1)}$:
\begin{equation}
    \mathbb{E}_{z \sim \mathcal{D}_\text{noisy}}[\kappa_n]
    - \mathbb{E}_{z \sim \mathcal{D}_\text{clean}}[\kappa_n]
    \;<\;
    \mathbb{E}_{z \sim \mathcal{D}_\text{noisy}}[\kappa^\dagger_n]
    - \mathbb{E}_{z \sim \mathcal{D}_\text{clean}}[\kappa^\dagger_n].
\end{equation}

\textbf{Memorization onset $T_m$.}
In the main paper (\S4.2.2) we define $T_m$ as the first checkpoint at which noisy-sample gradient norms begin to decrease.
The formal continuous analogue is:
\begin{equation}
    T_m := \inf\!\left\{t : \mathbb{E}_{z \sim \mathcal{D}_\text{noisy}}\!\left[\|\nabla_\theta \ell(z;\theta_t)\|\right] < \mathbb{E}_{z \sim \mathcal{D}_\text{noisy}}\!\left[\|\nabla_\theta \ell(z;\theta_0)\|\right]\right\}.
    \label{eq:Tm}
\end{equation}
\textbf{Assumptions.}
\begin{itemize}
    \item[\textbf{(A1)}] \textit{Monotone gradient decay for clean samples.}
        For $z \in \mathcal{D}_\text{clean}$, $I_{(n,t)}$ is non-increasing in $t$, converging toward zero (holds under $L$-smooth loss and a sufficiently small learning rate). The reciprocal $1/I_{(n,t)}$ is therefore non-decreasing.
    \item[\textbf{(A2)}] \textit{Bounded gradient for noisy samples within the proof horizon.}
        For $z \in \mathcal{D}_\text{noisy}$, $I_{(n,t)} \geq \delta > 0$ for all $t \leq T_m$, where $T_m$ is the memorization onset defined in Section~4.2.2 of the main paper.
        Beyond $T_m$, noisy gradients decay; the decreasing kernel discounts these late checkpoints.
        With $\lambda^* = T/T_m$, the kernel weight at $T_m+1$ relative to $t^*$ is $e^{-\lambda^*(T_m+1-t^*)/T}$, which decreases as $T_m \to T$, making the post-memorization contribution $S_2$ negligible relative to $S_1$.
    \item[\textbf{(A3)}] \textit{Post-memorization baseline inflation.}
        At $t = T-1 > T_m$, clean gradients are near zero (by A1) and noisy gradients have dropped under memorization.
        By symmetry near convergence, the noisy reciprocals exceed the clean reciprocals:
        $\Delta_{T-1} := \mathbb{E}_\text{noisy}[1/I_{(n,T-1)}] - \mathbb{E}_\text{clean}[1/I_{(n,T-1)}] > 0$.
    \item[\textbf{(A4)}] \textit{Positive temporal kernel.}
        $k(t) > 0$ for all $t$; decreasing so $k(t^*) > k(T-1)$ for any $t^* < T-1$.
\end{itemize}

\textbf{Proof.}

Define $\Delta_t := \mathbb{E}_\text{noisy}[1/I_{(n,t)}] - \mathbb{E}_\text{clean}[1/I_{(n,t)}]$, so
\begin{equation}
    \mathbb{E}_\text{noisy}[\kappa_n] - \mathbb{E}_\text{clean}[\kappa_n]
    = \sum_{t=0}^{T-1} k(t)\,\Delta_t.
\end{equation}

\textit{Step~1 --- Sign structure.}
For $1 \leq t \leq T_m$: (A1) and (A2) together give $\Delta_t \leq 0$, strictly negative at some $t^* \leq T_m$.
For $t > T_m$: $\Delta_t$ may be positive (memorization onset decreases noisy norms).

\textit{Step~2 --- Weighted sum.}
Split at $T_m$:
\begin{equation}
    S_1 = \sum_{t=1}^{T_m} k(t)\,\Delta_t \;<\; 0
    \qquad \text{and} \qquad
    S_2 = \sum_{t=T_m+1}^{T-1} k(t)\,\Delta_t \;\geq\; 0.
\end{equation}
Since the kernel is decreasing, $S_2 \leq k(T_m+1) \cdot \sum_{t > T_m} \Delta_t$, which is dominated by $|S_1|$ for sufficiently large $\lambda$.
Hence $\sum_t k(t)\,\Delta_t < 0$.

\textit{Step~3 --- Baseline gap.}
By (A3), $\mathbb{E}_\text{noisy}[\kappa^\dagger_n] - \mathbb{E}_\text{clean}[\kappa^\dagger_n] = \Delta_{T-1} > 0$.

\textit{Combining.}
\begin{equation}
    \mathbb{E}_\text{noisy}[\kappa_n] - \mathbb{E}_\text{clean}[\kappa_n]
    \;<\; 0
    \;<\; \Delta_{T-1}
    = \mathbb{E}_\text{noisy}[\kappa^\dagger_n] - \mathbb{E}_\text{clean}[\kappa^\dagger_n].
    \quad \square
\end{equation}

% ------------------------------------------------------------
\subsection{Distribution Matching Gap Bound (Theorem~1)}
\label{app:proof-gap}
% ------------------------------------------------------------

\textit{Cited in Section~4.4 of the main paper. Closes the chain from (RDM) back to (DD).}

\textbf{Goal.}
Let $f^*_\mathcal{S}$ be the model trained on $\mathcal{S}$ and $\mathcal{L}(\mathcal{S}) = \mathbb{E}_{z \sim P^w}[\ell(z; f^*_\mathcal{S})]$ be the downstream loss under $P^w$. Show:
\begin{equation}
    |\mathcal{L}(\tilde{\mathcal{D}}) - \mathcal{L}(\mathcal{D})|
    \;\leq\; 2L \cdot W_1(P^w, P_{\tilde{\mathcal{D}}}) + \beta_K.
\end{equation}

\textbf{Assumptions.}
\begin{itemize}
    \item[\textbf{(B1)}] $\ell(\cdot;\,f)$ is $L$-Lipschitz in the embedding metric $\|\phi(\cdot) - \phi(\cdot)\|$ for all $f$.
    \item[\textbf{(B2)}] The model class has uniform stability constant $\beta_K = O(1/K)$ \cite{bousquet2002stability}: $\sup_z |\ell(z; f^*_{\tilde{\mathcal{D}}}) - \ell(z; f^*_\mathcal{D})| \leq \beta_K$. Standard ERM algorithms (logistic regression, SVM) satisfy this for bounded loss; as $K \to \infty$, $\beta_K \to 0$.
    \item[\textbf{(B3)}] The ground metric for $W_1$ is $c(z,z') = \|\phi(z) - \phi(z')\|$.
\end{itemize}

\textbf{Step~1 --- Kantorovich--Rubinstein bound for a fixed model.}
By Kantorovich--Rubinstein duality \cite{villani2009optimal} under (B3), for any fixed $f$:
\begin{equation}
    \left|
        \mathbb{E}_{z \sim P^w}[\ell(z; f)]
        - \mathbb{E}_{z \sim P_{\tilde{\mathcal{D}}}}[\ell(z; f)]
    \right|
    \;\leq\; L \cdot W_1(P^w, P_{\tilde{\mathcal{D}}}).
\end{equation}

\textbf{Step~2 --- Triangle inequality decomposition.}
$\mathcal{L}(\tilde{\mathcal{D}})$ uses $f^*_{\tilde{\mathcal{D}}}$ while $\mathcal{L}(\mathcal{D})$ uses $f^*_\mathcal{D}$; both are evaluated under $P^w$ but with different models, so Step~1 does not apply directly.
Insert two intermediates and apply the triangle inequality twice:
\begin{align}
    |\mathcal{L}(\tilde{\mathcal{D}}) - \mathcal{L}(\mathcal{D})|
    &\leq
    \underbrace{
        \left|
            \mathbb{E}_{P^w}[\ell(z;f^*_{\tilde{\mathcal{D}}})]
            - \mathbb{E}_{P_{\tilde{\mathcal{D}}}}[\ell(z;f^*_{\tilde{\mathcal{D}}})]
        \right|
    }_{\text{(I) distribution shift, fixed }f^*_{\tilde{\mathcal{D}}}} \notag \\
    &\quad +
    \underbrace{
        \left|
            \mathbb{E}_{P_{\tilde{\mathcal{D}}}}[\ell(z;f^*_{\tilde{\mathcal{D}}})]
            - \mathbb{E}_{P_{\tilde{\mathcal{D}}}}[\ell(z;f^*_\mathcal{D})]
        \right|
    }_{\text{(II) model difference, fixed }P_{\tilde{\mathcal{D}}}} \notag \\
    &\quad +
    \underbrace{
        \left|
            \mathbb{E}_{P_{\tilde{\mathcal{D}}}}[\ell(z;f^*_\mathcal{D})]
            - \mathbb{E}_{P^w}[\ell(z;f^*_\mathcal{D})]
        \right|
    }_{\text{(III) distribution shift, fixed }f^*_\mathcal{D}}.
\end{align}
Terms~(I) and~(III) each satisfy Step~1 with $f = f^*_{\tilde{\mathcal{D}}}$ and $f = f^*_\mathcal{D}$ respectively, so each is bounded by $L \cdot W_1(P^w, P_{\tilde{\mathcal{D}}})$.
Term~(II) is bounded pointwise by (B2): $\text{(II)} \leq \beta_K = O(1/K)$.

\textbf{Combining.}
\begin{equation}
    |\mathcal{L}(\tilde{\mathcal{D}}) - \mathcal{L}(\mathcal{D})|
    \;\leq\; 2L \cdot W_1(P^w, P_{\tilde{\mathcal{D}}}) + \beta_K.
\end{equation}
Minimizing $W_1(P^w, P_{\tilde{\mathcal{D}}})$ via the OT objective directly minimizes the dominant term; $\beta_K \to 0$ as $K \to \infty$, a guarantee that k-means and greedy coreset methods lack. $\square$

\noindent\textit{References:} Kantorovich \& Rubinstein (1958); Villani (2009, Ch.~6).


% ============================================================
\section{Entropic Regularization for Tractability}
\label{app:method}
% ============================================================

This section extends the derivation in Section~4.4 of the main paper.
The discrete OT problem is defined over empirical measures $\mu = P^w = \sum_n w_n \delta_{z_n}$ (TAKE-weighted source) and $\nu = \frac{1}{|\mathcal{D}_\text{synth}|}\sum_c \delta_{\tilde{z}_c}$ (uniform target over the synthetic pool).

\textbf{Transportation cost (Kantorovich form).}
\begin{equation}
    \mathcal{C}(\gamma)
    = \int c(x,\hat{x})\,d\gamma(x,\hat{x}),
    \qquad c(x,\hat{x}) = \|\phi(x) - \phi(\hat{x})\|^2,
    \quad \gamma \in \Pi(\mu,\nu).
\end{equation}

\textbf{Entropic regularization.}
The bare Kantorovich problem is $O(N^3)$ complexity and non-differentiable.
Adding an entropic penalty yields a strictly convex, smooth objective:
\begin{equation}
    \min_{\gamma \in \Pi(\mu,\nu)}
    \langle C, \gamma \rangle
    + \varepsilon \cdot \mathrm{KL}(\gamma, \mu \otimes \nu),
    \qquad
    \mathrm{KL}(\gamma, \mu \otimes \nu)
    = \int \log\frac{d\gamma}{d(\mu \otimes \nu)}\,d\gamma,
\end{equation}
where $\mu \otimes \nu$ is the product (independence) reference coupling.
Smaller $\varepsilon$ recovers the exact OT solution; larger $\varepsilon$ smooths the plan toward the independent coupling, encouraging diversity in the selected prototypes.

\textbf{Optional prior coupling.}
If a task-specific prior $\alpha(x,\hat{x})$ is available (\eg, a topic-conditional measure), replace $\mu \otimes \nu$ with $\alpha$ to regularize toward domain knowledge.

\textbf{Sinkhorn--Knopp solver.}
The solution is computed by alternating row/column normalization on the kernel matrix $\mathbf{K}_{nc} = e^{-C_{nc}/\varepsilon}$, yielding $O(N \cdot |\mathcal{D}_\text{synth}|)$ per iteration.

\textbf{Prototype extraction.}
The distilled set $\tilde{\mathcal{D}}$ is formed by selecting the top-$K$ candidates ranked by received mass $\sum_n \gamma^*_{nc}$.


% ============================================================
\section{Implementation Details}
\label{app:expts}
% ============================================================

\textbf{Datasets.}
Table~\ref{tab:datasets} reports the task type and label set for each benchmark.

\begin{table}[ht]
    \caption{Dataset overview.}
    \label{tab:datasets}
    \centering
    \begin{tabular}{llrl}
        \toprule
        \textbf{Dataset} & \textbf{Task} & $N$ & \textbf{Labels} \\
        \midrule
        AG News   & Classification & 120,000 & World, Sports, Business, Sci/Tech \\
        IMDb      & Classification &  25,000 & Pos., Neg. \\
        SST-2     & Classification &  67,349 & Pos., Neg. \\
        MNLI-m    & NLI            & 392,702 & Entailment, Neutral, Contradiction \\
        QNLI      & NLI            & 104,743 & Entailment, Not Entailment \\
        QQP       & NLI            & 363,846 & Duplicate, Non-duplicate \\
        \bottomrule
    \end{tabular}
\end{table}

\textbf{Backbone models.}
Classification experiments use well-known classifiers such as Logistic Regression, TextCNN, TextRNN, and BERT.
NLI experiments use Siamese Logistic Regression and BERT.
These architectures are chosen as standard, reproducible benchmarks spanning the major NLP paradigms while remaining tractable at ablation study scale.

\begin{table}[ht]
    \caption{Backbone model details.}
    \label{tab:backbones}
    \centering
    \begin{tabular}{lll}
        \toprule
        \textbf{Model}        & \textbf{Architecture}        & \textbf{Tasks} \\
        \midrule
        Logistic              & Linear head on embeddings    & Classification \\
        TextCNN               & Convolutional encoder        & Classification \\
        TextRNN               & Recurrent encoder (GRU)      & Classification \\
        Siamese Logistic      & Siamese features \& logistic head & NLI \\
        BERT & Transformer encoder (110M) & Classification + NLI \\
        \bottomrule
    \end{tabular}
\end{table}

\textbf{Hyperparameters.}
All experiments run on a single NVIDIA V100 40\,GB GPU using PyTorch, PyTorch Lightning, HuggingFace Transformers, and the Python Optimal Transport (POT) library.
Influence scoring uses all-MiniLM-L6-v2~\cite{reimers2019sentence,wang2020minilm} as the frozen sentence encoder $\phi$.
We use $T = 20$ checkpoints, balancing trajectory coverage and compute.
The kernel decay is set to $\lambda^* = T / T_m$, concentrating mass on the pre-memorization phase as derived in Section~4.2.2.
For the Sinkhorn solver we use $\varepsilon = 0.05$ with 200 iterations, which produces soft plans with sufficient diversity and converges for all tested pool sizes.
The synthetic pool size is fixed at $|\mathcal{D}_\text{synth}| = 5{,}000$, chosen based on empirical saturation across all tested domains.
The language model generator is Gemma-3-270M, which is small, fast, and sufficient for in-domain fluency at this pool size.
